% Options for packages loaded elsewhere
\PassOptionsToPackage{unicode}{hyperref}
\PassOptionsToPackage{hyphens}{url}
\documentclass[
  11pt,
]{article}
\usepackage{xcolor}
\usepackage[margin=1in]{geometry}
\usepackage{amsmath,amssymb}
\setcounter{secnumdepth}{-\maxdimen} % remove section numbering
\usepackage{iftex}
\ifPDFTeX
  \usepackage[T1]{fontenc}
  \usepackage[utf8]{inputenc}
  \usepackage{textcomp} % provide euro and other symbols
\else % if luatex or xetex
  \usepackage{unicode-math} % this also loads fontspec
  \defaultfontfeatures{Scale=MatchLowercase}
  \defaultfontfeatures[\rmfamily]{Ligatures=TeX,Scale=1}
\fi
\usepackage{lmodern}
\ifPDFTeX\else
  % xetex/luatex font selection
\fi
% Use upquote if available, for straight quotes in verbatim environments
\IfFileExists{upquote.sty}{\usepackage{upquote}}{}
\IfFileExists{microtype.sty}{% use microtype if available
  \usepackage[]{microtype}
  \UseMicrotypeSet[protrusion]{basicmath} % disable protrusion for tt fonts
}{}
\makeatletter
\@ifundefined{KOMAClassName}{% if non-KOMA class
  \IfFileExists{parskip.sty}{%
    \usepackage{parskip}
  }{% else
    \setlength{\parindent}{0pt}
    \setlength{\parskip}{6pt plus 2pt minus 1pt}}
}{% if KOMA class
  \KOMAoptions{parskip=half}}
\makeatother
\setlength{\emergencystretch}{3em} % prevent overfull lines
\providecommand{\tightlist}{%
  \setlength{\itemsep}{0pt}\setlength{\parskip}{0pt}}
\usepackage{bookmark}
\IfFileExists{xurl.sty}{\usepackage{xurl}}{} % add URL line breaks if available
\urlstyle{same}
% fallback for those not using the hyperref driver hyperxmp:
\makeatletter
\@ifundefined{xmpquote}{\newcommand{\xmpquote}[1]{#1}}{}
\makeatother
\hypersetup{
  pdftitle={Simultaneous complement avoidance and explicit Thue--Morse barriers},
  hidelinks,
  pdfcreator={LaTeX via pandoc}}

\title{Simultaneous complement avoidance and explicit Thue--Morse
barriers}
\author{}
\date{20 September 2026 --- continuation for discussion}

\begin{document}
\maketitle

This continuation resolves two questions left in the earlier packet:
simultaneous avoidance of the two complementary Thue--Morse sequences,
and an explicit finite matching bound for the particular Thue--Morse
population. The first result holds for every non-eventually-periodic
binary sequence. These are hand proofs with finite computational
controls; independent review and priority checking remain. The earlier
Lean files do not certify these additions.

\section{1. One population avoids both complementary
targets}\label{one-population-avoids-both-complementary-targets}

For a binary sequence \(s\), write \(\bar s(k)=1-s(k)\) and put

\[r(2j)=s(j),\qquad r(2j+1)=1-s(j).\]

The population \(P_s\) has vertices \(\mathbb N\), birthdate \(v\) at
vertex \(v\), and, for every \(v\ge2\), precisely the incoming edges

\[v-1\longrightarrow v\text{ labelled }r(v),\qquad
v-2\longrightarrow v\text{ labelled }1-r(v).\]

There is no edge \(0\to1\). The roots are 0 and 1; every non-root has a
parent of each label; every vertex has at most two children. All
birthdate sublevel sets are finite. Thus this is an Alexander population
in the edge-labelled convention, with no restriction on a realizing
path's start.

\textbf{Theorem.} If \(s\) is not eventually periodic, \(P_s\) realizes
neither \(s\) nor \(\bar s\).

The original packet proves avoidance of \(s\). Here is a self-contained
joint argument. Let \(v_0,v_1,\ldots\) be a hypothetical matching path
and put \(d_k=v_k-2k\). Every increment of \(d_k\) is \(-1\) or 0.

For target \(s\), induction gives \(v_k\ge2k\): at \(v_k=2k\), the
increment-one edge, if present, has label \(r(2k+1)=1-s(k)\) and cannot
match. Hence \(d_k\ge0\) and eventually \(d_k=d\) is constant.

For target \(\bar s\), the corresponding induction gives \(v_k\ge2k-1\).
If \(v_k>2k-1\), any edge preserves the next lower bound. At equality,
an increment-one edge has label \(r(2k)=s(k)\) and cannot match. Thus
\(d_k\ge-1\) and again it stabilizes. In fact an infinite matching path
cannot visit \(d_k=-1\): both outgoing edges there have label \(s(k)\),
since the increment-two label is \(1-r(2k+1)=s(k)\). This also handles
that exceptional offset without invoking a shift of length zero.

On the constant-offset tail, every increment is two. For \(d\ge0\), put
\(q=\lfloor d/2\rfloor+1>0\). The matching equation is as follows:

For target \(s\): even \(d\) gives \(s(k)=1-s(k+q)\), and odd \(d\)
gives \(s(k)=s(k+q)\). For target \(\bar s\): even \(d\) gives
\(s(k)=s(k+q)\), and odd \(d\) gives \(s(k)=1-s(k+q)\).

Either equality makes \(s\) eventually periodic, with period \(q\) or
\(2q\). This proves the theorem. The boundary \(-1\), rather than 0, is
the only new ingredient needed for the complementary target.

For a periodic control of period \(p\), the construction correctly fails
to avoid either target: the increment-two paths starting at \(2p-1\) and
\(2p-2\) realize \(s\) and \(\bar s\), respectively. These controls
include \(p=1\) and the actual root convention.

\textbf{Consequence for the two-phase substitution gadget.} The earlier
exact relation is

\[\mu(x)\in L(Q(P))\quad\Longleftrightarrow\quad
x\in L(P)\text{ or }\bar x\in L(P),\qquad \mu(0)=01,\ \mu(1)=10.\]

Therefore \(Q(P_s)\) avoids \(\mu(s)\) and \(\mu(\bar s)\) whenever
\(s\) is aperiodic. For Thue--Morse \(t\), both \(t\) and \(\bar t\) are
fixed by \(\mu\), so \(P_t\) and every iterate \(Q^n(P_t)\) avoid both.
This settles the simultaneous-avoidance question that previously
prevented this particular use of the gadget. It does not construct an
intersection operation for arbitrary population languages.

\section{2. A short-run lemma for
Thue--Morse}\label{a-short-run-lemma-for-thuemorse}

Write \(t(k)=\operatorname{popcount}(k)\bmod2\), and, for \(q\ge1\),
define

\[c_q(k)=t(k)\mathbin{\mathrm{xor}}t(k+q),\qquad
\lambda(q)=2^{\nu_2(q)}.\]

\textbf{Lemma.} Every constant consecutive run of \(c_q\) has length at
most \(3\lambda(q)\), independently of where the run starts.

The identities \(t(2n)=t(n)\) and \(t(2n+1)=1-t(n)\) imply that \(t\)
has no three consecutive equal bits: any three positions include an
even/odd pair \((2n,2n+1)\).

First let \(q=2r+1\) be odd. Then

\[c_q(2n)=1\mathbin{\mathrm{xor}}t(n)\mathbin{\mathrm{xor}}t(n+r),\]
\[c_q(2n+1)=1\mathbin{\mathrm{xor}}t(n)\mathbin{\mathrm{xor}}t(n+r+1).\]

Four equal values starting at \(2n\) would force
\(t(n+r)=t(n+r+1)=t(n+r+2)\), by comparing the two aligned pairs. Four
equal values starting at \(2n+1\) would instead force
\(t(n)=t(n+1)=t(n+2)\), by comparing the pairs \((2n+1,2n+2)\) and
\((2n+3,2n+4)\). Both are impossible. Thus the bound is three for odd
\(q\).

Finally \(c_{2q}(2n)=c_{2q}(2n+1)=c_q(n)\). Each doubling of the shift
doubles each constant run. Removing all powers of two proves the lemma.
This elementary symbolic-sequence lemma is used here as a tool; no claim
that the lemma itself is new is intended.

\section{3. An explicit bound for every starting
vertex}\label{an-explicit-bound-for-every-starting-vertex}

For the Thue--Morse population \(P_t\), \(r=t\). Define

\[W(i)=\sum_{d=0}^{i}\lambda\!\left(\left\lfloor d/2\right\rfloor+1\right).\]

\textbf{Theorem.} A path starting at vertex \(i\) and matching a prefix
of \(t\) has at most \(i+3W(i)\) edges. A path matching a prefix of
\(\bar t\) has at most \(i+1+3W(i)\) edges. In particular neither target
can match

\[B(i)=i+2+3W(i)\]

edges from that vertex. This is a sufficient barrier, not a claimed
optimal one.

To prove it, separate increment-one edges from maximal consecutive runs
of increment-two edges. Each increment-one edge decreases \(d_k\) by
one; there are at most \(i\) such edges for target \(t\) and at most
\(i+1\) for target \(\bar t\). Each nonnegative offset \(d\) is visited
in at most one plateau. On its increment-two run, the equations above
make \(c_q(k)\) constant, where \(q=\lfloor d/2\rfloor+1\). The lemma
bounds this run by \(3\lambda(q)\). Summing over \(0\le d\le i\) proves
the stated bounds. Offset \(-1\) contributes no edges for the
complementary target.

The bound grows as \(O((i+1)\log(i+2))\). For an explicit estimate put
\(m=\lfloor i/2\rfloor+1\) and

\[S(m)=\sum_{q=1}^{m}\lambda(q)
=m+\sum_{j\ge1}2^{j-1}\left\lfloor m/2^j\right\rfloor
\le m\left(1+\tfrac12\lfloor\log_2m\rfloor\right).\]

Since \(W(i)\le2S(m)\),

\[B(i)\le i+2+6m+3m\lfloor\log_2m\rfloor.\]

For exact integer evaluation, \(S(0)=0\) and
\(S(m)=\lceil m/2\rceil+2S(\lfloor m/2\rfloor)\). Also
\(W(2h+1)=2S(h+1)\) and \(W(2h)=2S(h)+\lambda(h+1)\).

This replaces a qualitative termination argument for a \emph{fixed
target} with a computable a priori cutoff. It does not contradict the
earlier absence of a target-uniform bound at fixed start: an arbitrary
aperiodic target may begin with an arbitrarily long constant block.

\section{4. Checks, limitations and remaining
questions}\label{checks-limitations-and-remaining-questions}

The accompanying \texttt{evidence/44s-complement-barrier-checks.py}
enumerates finite paths directly from the graph's two incoming-edge
formulas, checks both boundary invariants on all short binary targets,
uses periodic positive controls, and compares actual Thue--Morse
first-empty times with \(B(i)\). It also checks the short-run lemma over
finite ranges. The JSON records the precise ranges and outcomes; finite
checks are not the infinite proofs.

The logarithmic factor may be avoidable: the proof separately charges
every offset's longest possible run even when those runs cannot be
concatenated along one matching path. A sharp asymptotic or exact
barrier remains open locally. \textbf{September 21 continuation:}
\texttt{countable-avoidance.md} now settles simultaneous avoidance for
arbitrary countable aperiodic families, including finite collections, at
optimal degree for every finite alphabet. Other embedding/homomorphism
versions of universal avoiding populations and more informative ordinal
invariants remain outside these results.

Source: Samuel A. Alexander, \emph{Biologically unavoidable sequences},
Electronic Journal of Combinatorics \textbf{20}(1) (2013), P31,
\href{https://doi.org/10.37236/3035}{doi:10.37236/3035}. The population
and earlier classification proof are in this packet's
\texttt{results.md}. The claims here are local extensions of that
construction, not attributed to Alexander.

\end{document}
